\usepackage{MinionPro} 
\usepackage[margin=1.5in]{geometry}
% \usepackage{tgtermes}
% \usepackage{amssymb}
% \usepackage{newtxmath}

\usepackage{lettrine}
\usepackage{mathrsfs}
\usepackage{comment}
\usepackage[usenames,dvipsnames]{color}
\usepackage[normalem]{ulem}
\usepackage{url}
\usepackage{graphicx}
\usepackage{float, caption}
\usepackage{enumerate}
\usepackage{hyperref} 
\usepackage[dvipsnames]{xcolor}
\usepackage{tikz}
\usepackage{quiver}
\usepackage{anyfontsize}
\usepackage{epigraph}
\setlength\epigraphrule{0pt}

\usepackage{etoolbox}
\usepackage[shortlabels]{enumitem}


% Make \section Large and bold, centered
\patchcmd{\section}{\scshape}{\Large \scshape}{}{}
% Changes the word "Contents" to "Table of Contents" and makes it Large
\renewcommand{\contentsname}{\Large Contents}


\makeatletter
% Section: level 1, 0pt vertical space, 0pt indent, bold
\renewcommand{\l@section}{\@tocline{1}{3pt}{0pt}{}{\normalfont }}

% Subsection: level 2, 0pt vertical space, 2.5pc indent, normal font
\renewcommand{\l@subsection}{\@tocline{2}{0pt}{2.5pc}{}{\normalfont}}
\makeatother


\usetikzlibrary{decorations.pathreplacing, calc}
\tikzcdset{
  diagrams={>={Straight Barb[scale=0.5]}}
}

\tikzset{
  COSIMPLICIALaltstackar/.style={decorate, decoration={show path construction,
    lineto code={
      \path (\tikzinputsegmentfirst); \pgfgetlastxy{\xstart}{\ystart}
      \path (\tikzinputsegmentlast); \pgfgetlastxy{\xend}{\yend}
      \path ($(0,0)!3pt!(\ystart-\yend,\xend-\xstart)$); \pgfgetlastxy{\xperp}{\yperp}
      \foreach \n[evaluate=\n as \k using .5*#1-\n+.5] in {1,...,#1}{
        \ifodd\n{\draw[->, shorten <=2pt, shift={($\k*(\xperp,\yperp)$)}](\xstart,\ystart)--(\xend,\yend);}
        \else{\draw[<-, shorten >=2pt, shift={($\k*(\xperp,\yperp)$)}](\xstart,\ystart)--(\xend,\yend);}\fi
      }
    }
  }}, COSIMPLICIALaltstackar/.default={1}
}


\tikzset{
  SIMPLICIALaltstackar/.style={decorate, decoration={show path construction,
    lineto code={
      \path (\tikzinputsegmentfirst); \pgfgetlastxy{\xstart}{\ystart}
      \path (\tikzinputsegmentlast); \pgfgetlastxy{\xend}{\yend}
      \path ($(0,0)!3pt!(\ystart-\yend,\xend-\xstart)$); \pgfgetlastxy{\xperp}{\yperp}
      \foreach \n[evaluate=\n as \k using .5*#1-\n+.5] in {1,...,#1}{
        \ifodd\n{\draw[<-, shorten >=2pt, shift={($\k*(\xperp,\yperp)$)}](\xstart,\ystart)--(\xend,\yend);}
        \else{\draw[->, shorten <=2pt, shift={($\k*(\xperp,\yperp)$)}](\xstart,\ystart)--(\xend,\yend);}\fi
      }
    }
  }}, SIMPLICIALaltstackar/.default={1}
}


\hypersetup{%
  bookmarksnumbered=true,%
  colorlinks=true,
  allcolors= RoyalPurple
}

% 2. cleveref MUST be loaded after hyperref, but before \newtheorems
\usepackage[capitalize,nameinlink]{cleveref} 

\let\fullref\autoref
\def\makeautorefname#1#2{\expandafter\def\csname#1autorefname\endcsname{#2}}
\def\equationautorefname~#1\null{(#1)\null}
\makeautorefname{footnote}{footnote}%
\makeautorefname{item}{item}%
\makeautorefname{figure}{Figure}%
\makeautorefname{table}{Table}%
\makeautorefname{part}{Part}%
\makeautorefname{appendix}{Appendix}%
\makeautorefname{chapter}{Chapter}%
\makeautorefname{section}{Section}%
\makeautorefname{subsection}{Section}%
\makeautorefname{subsubsection}{Section}%
\makeautorefname{theorem}{Theorem}%
\makeautorefname{thm}{Theorem}%
\makeautorefname{cor}{Corollary}%
\makeautorefname{lem}{Lemma}%
\makeautorefname{prop}{Proposition}%
\makeautorefname{pro}{Property}
\makeautorefname{conj}{Conjecture}%
\makeautorefname{defn}{Definition}%
\makeautorefname{notn}{Notation}
\makeautorefname{notns}{Notations}
\makeautorefname{rem}{Remark}%
\makeautorefname{quest}{Question}%
\makeautorefname{exmp}{Example}%
\makeautorefname{ax}{Axiom}%
\makeautorefname{claim}{Claim}%
\makeautorefname{ass}{Assumption}%
\makeautorefname{asss}{Assumptions}%
\makeautorefname{con}{Construction}%
\makeautorefname{warn}{Warning}%
\makeautorefname{obs}{Observation}%
\makeautorefname{conv}{Convention}%

% 3. Defining environments with standard shared counter syntax [thm] 
% This natively handles what the \makeatletter hack was trying to do.
%theoremstyle{plain} --- default
\newtheorem{thm}{Theorem}[section]
\newtheorem{cor}[thm]{Corollary}
\newtheorem{prop}[thm]{Proposition}
\newtheorem{lem}[thm]{Lemma}
\newtheorem*{exer}{Exercise}
\newtheorem{conj}[thm]{Conjecture}

\theoremstyle{definition}
\newtheorem{defn}[thm]{Definition}
\newtheorem{ass}[thm]{Assumption}
\newtheorem{asss}[thm]{Assumptions}
\newtheorem{ax}[thm]{Axiom}
\newtheorem{con}[thm]{Construction}
\newtheorem{exmp}[thm]{Example}
\newtheorem{notn}[thm]{Notation}
\newtheorem{notns}[thm]{Notations}
\newtheorem{pro}[thm]{Property}
\newtheorem{quest}[thm]{Question}
\newtheorem{rem}[thm]{Remark}
\newtheorem{warn}[thm]{Warning}
\newtheorem{sch}[thm]{Scholium}
\newtheorem{obs}[thm]{Observation}
\newtheorem{conv}[thm]{Convention}
\newtheorem{claim}[thm]{Claim}

\newenvironment{soln}{\begin{proof}[Solution]}{\end{proof}}

% Reset equations to number by section
\numberwithin{equation}{section}

% ===== Math Operators =====
\DeclareMathOperator{\ex}{ex}
\DeclareMathOperator{\ch}{ch}
\DeclareMathOperator*{\argmax}{arg\,max}
\DeclareMathOperator*{\argmin}{arg\,min}
\DeclareMathOperator{\epi}{epi}
\DeclareMathOperator{\rec}{rec}
\DeclareMathOperator{\csh}{csh}
\DeclareMathOperator{\sect}{Sect}
\DeclareMathOperator{\adj}{adj}
\DeclareMathOperator{\vertices}{vert}
\DeclareMathOperator{\cl}{cl}
\DeclareMathOperator{\im}{im}
\DeclareMathOperator{\atan2}{atan2}
\DeclareMathOperator{\arc}{arc}
\DeclareMathOperator{\edge}{edge}
\DeclareMathOperator{\tensor}{\otimes}
\DeclareMathOperator{\tot}{Tot}
\DeclareMathOperator{\dir}{dir}
\DeclareMathOperator{\sweep}{sweep}
\DeclareMathOperator{\comet}{Comet}
\DeclareMathOperator{\nor}{nor}
\DeclareMathOperator{\Map}{Map}
\DeclareMathOperator{\Cyl}{Cyl}
\DeclareMathOperator{\hofib}{hofib}
\DeclareMathOperator{\colim}{colim}
\DeclareMathOperator{\id}{id}
\DeclareMathOperator{\conf}{Conf}
\DeclareMathOperator{\Inj}{Inj}
\DeclareMathOperator{\Intr}{int}
\DeclareMathOperator{\Vect}{Vect}
\DeclareMathOperator{\proj}{proj}
\DeclareMathOperator{\AW}{\mathcal{AW}}
\DeclareMathOperator{\EZ}{\mathcal{EZ}}
\DeclareMathOperator{\Sing}{Sing}
\DeclareMathOperator{\coker}{coker}
\DeclareMathOperator{\Cobar}{\catname{Cobar}}
\DeclareMathOperator{\Barr}{\catname{Bar}}
\DeclareMathOperator{\der}{Der}
\DeclareMathOperator{\Hom}{Hom}
\DeclareMathOperator{\gr}{gr}
\DeclareMathOperator{\GL}{GL}
\DeclareMathOperator{\Sym}{Sym}
\DeclareMathOperator{\Mod}{\catname{Mod}}
\DeclareMathOperator{\Ch}{\catname{Ch}}

% ===== Custom Commands =====
\newcommand{\after}{$\circ$}
\newcommand{\N}{\mathbf{N}}
\newcommand{\Z}{\mathbf{Z}}
\newcommand{\Q}{\mathbf{Q}}
\newcommand{\R}{\mathbf{R}}
\newcommand{\C}{\mathbf{C}}
\newcommand{\RP}{\mathbf{R}P}
\newcommand{\gp}{\mathrm{gp}}
\newcommand{\mfm}{\mathfrak{m}}
\newcommand{\mfp}{\mathfrak{p}}
\newcommand{\mco}{\mathcal{O}}




\newcommand{\CP}{\mathbf{C}P}
\newcommand{\mb}[1]{\mathbf{#1}}
\newcommand{\lb}{\langle}
\newcommand{\rb}{\rangle}
\newcommand{\F}{\mathbf{F}}
\newcommand{\HH}{\mathbf{H}}
\newcommand{\T}{\mathcal{T}}
\newcommand{\pset}{\mathscr{P}}
\newcommand{\catname}{\mathsf}
\newcommand{\p}{^{\prime}}
\newcommand{\pp}{^{\prime\prime}}
\newcommand{\iso}{\cong}
\newcommand{\eps}{\varepsilon}
\newcommand{\dx}{\,dx}
\newcommand{\dy}{\,dy}
\newcommand{\dt}{\,dt}
\newcommand{\dz}{\,dz}
\newcommand{\dvol}{d\!\vol}
\newcommand{\ceq}{\coloneqq}
\newcommand{\tto}{\rightrightarrows}
\newcommand{\nat}{\Rightarrow}
\newcommand{\contains}{\supseteq}
\newcommand{\onto}{\twoheadrightarrow}
\newcommand{\into}{\hookrightarrow}
\newcommand{\bigtensor}{\bigotimes}
\newcommand{\const}{\mathrm{const}}
\newcommand{\CCobar}{\mathfrak{C}\catname{obar}}
\newcommand{\res}{\mathrm{res}}
\newcommand{\shom}{\mathcal{H}om}


\renewcommand{\qedsymbol}{$\square$}
\newcommand{\psh}[2]{\operatorname{PSh}(\mathcal{#1}, \mathcal{#2})}
\makeatother